% =============================================================================

% Dokumentacja projektu: Hybrydowy model LSTM + Transformer

% Prognozowanie kierunku kursu USD/PLN

%

% Kompilacja (w katalogu Kod/):

%   pdflatex dokumentacja.tex

%   pdflatex dokumentacja.tex   % drugi raz — spis treści

%

% Wymagane pakiety: babel, amsmath, amssymb, graphicx, hyperref,

%                   booktabs, listings, float, geometry

% =============================================================================



\documentclass[12pt,a4paper]{report}



\usepackage[utf8]{inputenc}

\usepackage[T1]{fontenc}

\usepackage[polish]{babel}

\usepackage{amsmath, amssymb}

\usepackage{graphicx}

\usepackage{hyperref}

\usepackage{booktabs}

\usepackage{listings}

\usepackage{float}

\usepackage{geometry}

\usepackage{array}

\usepackage{longtable}



\geometry{margin=2.5cm}

\hypersetup{colorlinks=true, linkcolor=black, urlcolor=blue}



\lstset{

    language=Python,

    basicstyle=\ttfamily\small,

    keywordstyle=\bfseries,

    commentstyle=\itshape,

    numbers=left,

    numberstyle=\tiny,

    stepnumber=1,

    numbersep=8pt,

    frame=single,

    breaklines=true,

    showstringspaces=false,

    tabsize=4,

}



% --- Uzupełnij imię, nazwisko i numer grupy ---
\newcommand{\autorDokumentacji}{Imię Nazwisko}
\newcommand{\grupaStudencka}{Grupa X}

\begin{document}

\begin{titlepage}
    \centering
    \vspace*{1.5cm}

    {\LARGE\bfseries Prognozowanie kierunku kursu USD/PLN\par}
    \vspace{0.8cm}
    {\Large Hybrydowy model LSTM\,+\,Transformer\par}
    \vspace{1.2cm}
    {\large Projekt z przedmiotu Sztuczna inteligencja\par}
    \vspace{0.4cm}
    {\Large\bfseries Dokumentacja\par}

    \vfill

    {\large \autorDokumentacji\par}
    \vspace{0.6cm}
    {\normalsize INFORMATYKA STOSOWANA\par}
    {\normalsize Semestr IV $\diamond$ \grupaStudencka\par}

    \vspace{1.5cm}

    {\large Uniwersytet Technologiczno-Przyrodniczy\par}
    {\large im.\ Jana i Jędrzeja Śniadeckich w Bydgoszczy\par}
    \vspace{0.4cm}
    {\normalsize Wydział Telekomunikacji, Informatyki i Elektrotechniki\par}

    \vspace{1.5cm}
    {\large Bydgoszcz, 2026\par}
\end{titlepage}

\tableofcontents

\newpage



% ============================================================================

\chapter{Wstęp}

% ============================================================================



% =============================================================================
% Wstęp — sekcja do wklejenia w istniejący dokument LaTeX
%
% W dokumencie nadrzędnym użyj jednego z poniższych wariantów:
%   \chapter{Wstęp}   % klasa report / book
%   \section{Wstęp}   % klasa article
%
% Wymagane w preambule dokumentu nadrzędnego:
%   \usepackage[polish]{babel}
%   \usepackage{amsmath, amssymb}
%
% Dołączenie: % =============================================================================
% Wstęp — sekcja do wklejenia w istniejący dokument LaTeX
%
% W dokumencie nadrzędnym użyj jednego z poniższych wariantów:
%   \chapter{Wstęp}   % klasa report / book
%   \section{Wstęp}   % klasa article
%
% Wymagane w preambule dokumentu nadrzędnego:
%   \usepackage[polish]{babel}
%   \usepackage{amsmath, amssymb}
%
% Dołączenie: % =============================================================================
% Wstęp — sekcja do wklejenia w istniejący dokument LaTeX
%
% W dokumencie nadrzędnym użyj jednego z poniższych wariantów:
%   \chapter{Wstęp}   % klasa report / book
%   \section{Wstęp}   % klasa article
%
% Wymagane w preambule dokumentu nadrzędnego:
%   \usepackage[polish]{babel}
%   \usepackage{amsmath, amssymb}
%
% Dołączenie: \input{wstep}
% =============================================================================

% --- Odkomentuj odpowiednią linię lub umieść treść poniżej pod własnym nagłówkiem ---

% \chapter{Wstęp}
% \section{Wstęp}

Rynek walutowy stanowi jeden z najbardziej płynnych i globalnie zintegrowanych segmentów rynków finansowych.
Kurs wymiany dolara amerykańskiego względem złotego polskiego (USD/PLN) ma bezpośrednie znaczenie dla polskiej gospodarki:
wpływa na koszty importu, wartość zobowiązań zadłużonych w walucie obcej oraz poziom inflacji importowanej.
Instytucje finansowe, przedsiębiorstwa i inwestorzy indywidualni poszukują narzędzi wspierających podejmowanie decyzji na tym rynku,
jednak prognozowanie kursów walutowych pozostaje zadaniem wyjątkowo trudnym.

Szeregi czasowe notowań finansowych charakteryzują się wysokim poziomem szumu, nieliniowością zależności
oraz okresową zmianą reżimów rynkowych. Klasyczne modele ekonometryczne, oparte na założeniach stacjonarności
i liniowości, często nie wystarczają do uchwycenia złożonych wzorców zachowania cen.
W ostatnich latach coraz większą popularnością cieszą się metody uczenia głębokiego, w szczególności sieci rekurencyjne
typu LSTM (Long Short-Term Memory) oraz architektury oparte na mechanizmie uwagi (attention), takie jak Transformer.
Pozwalają one na modelowanie wielowymiarowych sekwencji czasowych i łączenie informacji z różnych horyzontów temporalnych.

\section*{Problem badawczy}

Niniejsza praca podejmuje problem prognozowania kierunku zmiany kursu USD/PLN w horyzoncie 14~dni roboczych.
Główne pytanie badawcze brzmi: \emph{czy na podstawie 30-dniowej historii cech technicznych
i makroekonomicznych można skutecznie przewidzieć, czy kurs wzrośnie, czy spadnie w ciągu kolejnych 14~dni?}

Punktem wyjścia jest skumulowany zwrot procentowy kursu zamknięcia w horyzoncie prognozy.
Ze względu na praktyczną interpretację wyników (sygnały kupna i sprzedaży) zadanie sprowadzono
do klasyfikacji binarnej: klasa~1 oznacza wzrost kursu, klasa~0 oznacza spadek.
Etykieta dla dnia $t$ definiowana jest jako:
\begin{equation}
y_t =
\begin{cases}
1 & \text{gdy } \dfrac{P_{t+14}}{P_t} - 1 > 0, \\[6pt]
0 & \text{gdy } \dfrac{P_{t+14}}{P_t} - 1 \le 0.
\end{cases}
\label{eq:target}
\end{equation}
gdzie $P_t$ oznacza kurs zamknięcia pary USD/PLN w dniu~$t$.

Skumulowany zwrot w procentach wyraża się natomiast wzorem:
\begin{equation}
    r_t = \left(\frac{P_{t+14}}{P_t} - 1\right) \cdot 100\,\%.
    \label{eq:return}
\end{equation}

\section*{Cel pracy}

Celem pracy jest zaprojektowanie, implementacja i ocena hybrydowego modelu głębokiego łączącego warstwę LSTM
z blokiem Transformer w zadaniu klasyfikacji kierunku kursu USD/PLN.
W ramach realizacji celu:
\begin{itemize}
    \item przygotowano wielowymiarowy zestaw cech obejmujący wskaźniki analizy technicznej
          oraz stopy procentowe banków centralnych (Fed, EBC, NBP),
    \item zbudowano i wytrenowano model na historycznych danych dziennych,
    \item oceniono skuteczność predykcji za pomocą metryk klasyfikacji
          (accuracy, precision, recall, F1) oraz porównano wyniki z prostymi benchmarkami
          (strategia ,,zawsze wzrost'' i ,,zawsze spadek'').
\end{itemize}

\section*{Opis metody}

\subsection*{Dane wejściowe}

Analizie poddano dzienne notowania z pliku \texttt{dane\_ekonomiczne.csv}, którego źródłem jest Yahoo Finance oraz World Bank Data.
Dane obejmują m.in.\ ceny OHLC pary USD/PLN, skorelowane pary walutowe (EUR/PLN, EUR/USD)
oraz stopy procentowe banków centralnych Stanów Zjednoczonych, strefy euro i Polski.
Z zestawu usunięto obserwacje przypadające na weekendy, aby zachować ciągłość danych sesyjnych.
Przed przekazaniem do modelu wszystkie cechy numeryczne standaryzowano za pomocą \texttt{StandardScaler},
dopasowanego wyłącznie na zbiorze treningowym.

\subsection*{Inżynieria cech}

Dla każdego dnia obliczono łącznie 24~zmienne wejściowe, w tym:
\begin{itemize}
    \item dzienne stopy zwrotu kursów zamknięcia (USD/PLN, EUR/PLN, EUR/USD),
    \item wskaźniki momentum za okresy 3, 5 i 10~dni,
    \item RSI (14~okresów), ATR (14~okresów), pozycję i szerokość pasm Bollingera (okno 20~dni),
    \item wskaźniki MACD wraz z linią sygnałową i histogramem,
    \item zmienność kroczącą, $z$-score (okno 20~dni), dzienny zakres High/Low oraz stosunek High/Low,
    \item stopy procentowe: US\_FED, EA\_ECB, PL\_NBP,
    \item opóźnione kierunki zmian ceny (lagi 1, 2, 3 i 5~dni).
\end{itemize}

\subsection*{Architektura modelu}

Model przyjmuje na wejściu sekwencję $X_t \in \mathbb{R}^{30 \times 24}$,
tj.\ 30~kolejnych dni obserwacji po 24~cechy (okno przesuwne).
Architektura hybrydowa składa się z następujących etapów:
\begin{enumerate}
    \item warstwa LSTM z 128~jednostkami ukrytymi (\texttt{return\_sequences=True}),
          wychwytująca lokalne zależności czasowe w sekwencji,
    \item warstwa \texttt{Dropout} ($p = 0{,}3$) redukująca przeuczenie,
    \item kodowanie pozycyjne (\emph{Positional Encoding}) nadające informację o położeniu obserwacji w oknie,
    \item blok Transformer z 8~głowicami uwagi (\emph{Multi-Head Attention}),
          $d_{\text{model}} = 128$ i warstwą sprzężeń zwrotnych o wymiarze 256,
          modelujący długoterminowe zależności w obrębie 30-dniowego okna,
    \item agregacja \texttt{GlobalAveragePooling1D}, po której następują warstwy gęste
          (64 i 16~neuronów, aktywacja ReLU, regularyzacja L2) oraz warstwa wyjściowa
          z aktywacją sigmoid (jedna neurona, klasyfikacja binarna).
\end{enumerate}

Połączenie LSTM i Transformera uzasadnia fakt, że sieć rekurencyjna efektywnie koduje krótkoterminowe wzorce cenowe,
natomiast mechanizm self-attention pozwala ważyć znaczenie poszczególnych dni w historycznym oknie,
co jest istotne przy analizie wielotygodniowych trendów.

\subsection*{Trenowanie i ewaluacja}

Dane podzielono chronologicznie: 80\,\% stanowi zbiór treningowy, 20\,\%~--- testowy.
Z części treningowej wydzielono 10\,\% obserwacji na walidację.
Ze względu na naturę szeregu czasowego podczas uczenia nie stosowano losowego mieszania próbek (\texttt{shuffle=False}).

Model trenowano algorytmem Adam z początkową szybkością uczenia $5 \times 10^{-5}$,
funkcją straty \texttt{binary\_crossentropy} oraz wagami klas (\texttt{compute\_class\_weight})
kompensującymi niezbalansowanie etykiet.
Zastosowano wczesne zatrzymanie (\emph{early stopping}) monitorujące dokładność walidacyjną
oraz redukcję szybkości uczenia przy stagnacji straty.
Predykcje ciągłe z warstwy sigmoid konwertowano na klasy za pomocą progu $0{,}5$.

Jakość modelu oceniano metrykami: accuracy, precision, recall, F1 oraz macierzą pomyłek.
Wyniki porównywano z naiwnymi strategiami bazowymi przewidującymi zawsze wzrost lub zawsze spadek kursu.

\section*{Zakres i ograniczenia}

Praca koncentruje się wyłącznie na prognozie \emph{kierunku} zmiany ceny, a nie na estymacji jej poziomu
ani wielkości zwrotu. Nie uwzględniono kosztów transakcyjnych, spreadu bid-ask ani dźwigni finansowej,
co ogranicza bezpośrednią aplikowalność wyników w realnym handlu algorytmicznym.

Model wytrenowano na danych historycznych z okresu od 2020~roku; jego skuteczność w przyszłości
nie jest gwarantowana ze względu na ryzyko przeuczenia oraz możliwość zmiany reżimu rynkowego
(np.\ interwencje banków centralnych, kryzysy geopolityczne). Wyniki należy traktować jako eksperyment badawczy,
a nie rekomendację inwestycyjną.

\section*{Struktura pracy}

% --- Dostosuj nazwy rozdziałów do własnego spisu treści ---

Praca składa się z następujących rozdziałów:
\begin{enumerate}
    \item Wstęp --- motywacja, problem badawczy, cel i zakres pracy.
    \item Przegląd literatury i podstawy teoretyczne --- szeregi czasowe finansowe,
          sieci LSTM, architektura Transformer oraz wybrane wskaźniki analizy technicznej.
    \item Przygotowanie danych i inżynieria cech --- opis źródeł danych, preprocessingu
          i konstrukcji 24~zmiennych wejściowych.
    \item Architektura i trenowanie modelu --- szczegółowy opis hybrydy LSTM\,+\,Transformer,
          hiperparametrów i procedury walidacji.
    \item Wyniki i analiza --- metryki klasyfikacji, macierz pomyłek, porównanie z benchmarkami
          oraz wizualizacja sygnałów kupna i sprzedaży.
    \item Podsumowanie i wnioski --- ocena osiągniętych rezultatów i kierunki dalszych badań.
\end{enumerate}

% =============================================================================

% --- Odkomentuj odpowiednią linię lub umieść treść poniżej pod własnym nagłówkiem ---

% \chapter{Wstęp}
% \section{Wstęp}

Rynek walutowy stanowi jeden z najbardziej płynnych i globalnie zintegrowanych segmentów rynków finansowych.
Kurs wymiany dolara amerykańskiego względem złotego polskiego (USD/PLN) ma bezpośrednie znaczenie dla polskiej gospodarki:
wpływa na koszty importu, wartość zobowiązań zadłużonych w walucie obcej oraz poziom inflacji importowanej.
Instytucje finansowe, przedsiębiorstwa i inwestorzy indywidualni poszukują narzędzi wspierających podejmowanie decyzji na tym rynku,
jednak prognozowanie kursów walutowych pozostaje zadaniem wyjątkowo trudnym.

Szeregi czasowe notowań finansowych charakteryzują się wysokim poziomem szumu, nieliniowością zależności
oraz okresową zmianą reżimów rynkowych. Klasyczne modele ekonometryczne, oparte na założeniach stacjonarności
i liniowości, często nie wystarczają do uchwycenia złożonych wzorców zachowania cen.
W ostatnich latach coraz większą popularnością cieszą się metody uczenia głębokiego, w szczególności sieci rekurencyjne
typu LSTM (Long Short-Term Memory) oraz architektury oparte na mechanizmie uwagi (attention), takie jak Transformer.
Pozwalają one na modelowanie wielowymiarowych sekwencji czasowych i łączenie informacji z różnych horyzontów temporalnych.

\section*{Problem badawczy}

Niniejsza praca podejmuje problem prognozowania kierunku zmiany kursu USD/PLN w horyzoncie 14~dni roboczych.
Główne pytanie badawcze brzmi: \emph{czy na podstawie 30-dniowej historii cech technicznych
i makroekonomicznych można skutecznie przewidzieć, czy kurs wzrośnie, czy spadnie w ciągu kolejnych 14~dni?}

Punktem wyjścia jest skumulowany zwrot procentowy kursu zamknięcia w horyzoncie prognozy.
Ze względu na praktyczną interpretację wyników (sygnały kupna i sprzedaży) zadanie sprowadzono
do klasyfikacji binarnej: klasa~1 oznacza wzrost kursu, klasa~0 oznacza spadek.
Etykieta dla dnia $t$ definiowana jest jako:
\begin{equation}
y_t =
\begin{cases}
1 & \text{gdy } \dfrac{P_{t+14}}{P_t} - 1 > 0, \\[6pt]
0 & \text{gdy } \dfrac{P_{t+14}}{P_t} - 1 \le 0.
\end{cases}
\label{eq:target}
\end{equation}
gdzie $P_t$ oznacza kurs zamknięcia pary USD/PLN w dniu~$t$.

Skumulowany zwrot w procentach wyraża się natomiast wzorem:
\begin{equation}
    r_t = \left(\frac{P_{t+14}}{P_t} - 1\right) \cdot 100\,\%.
    \label{eq:return}
\end{equation}

\section*{Cel pracy}

Celem pracy jest zaprojektowanie, implementacja i ocena hybrydowego modelu głębokiego łączącego warstwę LSTM
z blokiem Transformer w zadaniu klasyfikacji kierunku kursu USD/PLN.
W ramach realizacji celu:
\begin{itemize}
    \item przygotowano wielowymiarowy zestaw cech obejmujący wskaźniki analizy technicznej
          oraz stopy procentowe banków centralnych (Fed, EBC, NBP),
    \item zbudowano i wytrenowano model na historycznych danych dziennych,
    \item oceniono skuteczność predykcji za pomocą metryk klasyfikacji
          (accuracy, precision, recall, F1) oraz porównano wyniki z prostymi benchmarkami
          (strategia ,,zawsze wzrost'' i ,,zawsze spadek'').
\end{itemize}

\section*{Opis metody}

\subsection*{Dane wejściowe}

Analizie poddano dzienne notowania z pliku \texttt{dane\_ekonomiczne.csv}, którego źródłem jest Yahoo Finance oraz World Bank Data.
Dane obejmują m.in.\ ceny OHLC pary USD/PLN, skorelowane pary walutowe (EUR/PLN, EUR/USD)
oraz stopy procentowe banków centralnych Stanów Zjednoczonych, strefy euro i Polski.
Z zestawu usunięto obserwacje przypadające na weekendy, aby zachować ciągłość danych sesyjnych.
Przed przekazaniem do modelu wszystkie cechy numeryczne standaryzowano za pomocą \texttt{StandardScaler},
dopasowanego wyłącznie na zbiorze treningowym.

\subsection*{Inżynieria cech}

Dla każdego dnia obliczono łącznie 24~zmienne wejściowe, w tym:
\begin{itemize}
    \item dzienne stopy zwrotu kursów zamknięcia (USD/PLN, EUR/PLN, EUR/USD),
    \item wskaźniki momentum za okresy 3, 5 i 10~dni,
    \item RSI (14~okresów), ATR (14~okresów), pozycję i szerokość pasm Bollingera (okno 20~dni),
    \item wskaźniki MACD wraz z linią sygnałową i histogramem,
    \item zmienność kroczącą, $z$-score (okno 20~dni), dzienny zakres High/Low oraz stosunek High/Low,
    \item stopy procentowe: US\_FED, EA\_ECB, PL\_NBP,
    \item opóźnione kierunki zmian ceny (lagi 1, 2, 3 i 5~dni).
\end{itemize}

\subsection*{Architektura modelu}

Model przyjmuje na wejściu sekwencję $X_t \in \mathbb{R}^{30 \times 24}$,
tj.\ 30~kolejnych dni obserwacji po 24~cechy (okno przesuwne).
Architektura hybrydowa składa się z następujących etapów:
\begin{enumerate}
    \item warstwa LSTM z 128~jednostkami ukrytymi (\texttt{return\_sequences=True}),
          wychwytująca lokalne zależności czasowe w sekwencji,
    \item warstwa \texttt{Dropout} ($p = 0{,}3$) redukująca przeuczenie,
    \item kodowanie pozycyjne (\emph{Positional Encoding}) nadające informację o położeniu obserwacji w oknie,
    \item blok Transformer z 8~głowicami uwagi (\emph{Multi-Head Attention}),
          $d_{\text{model}} = 128$ i warstwą sprzężeń zwrotnych o wymiarze 256,
          modelujący długoterminowe zależności w obrębie 30-dniowego okna,
    \item agregacja \texttt{GlobalAveragePooling1D}, po której następują warstwy gęste
          (64 i 16~neuronów, aktywacja ReLU, regularyzacja L2) oraz warstwa wyjściowa
          z aktywacją sigmoid (jedna neurona, klasyfikacja binarna).
\end{enumerate}

Połączenie LSTM i Transformera uzasadnia fakt, że sieć rekurencyjna efektywnie koduje krótkoterminowe wzorce cenowe,
natomiast mechanizm self-attention pozwala ważyć znaczenie poszczególnych dni w historycznym oknie,
co jest istotne przy analizie wielotygodniowych trendów.

\subsection*{Trenowanie i ewaluacja}

Dane podzielono chronologicznie: 80\,\% stanowi zbiór treningowy, 20\,\%~--- testowy.
Z części treningowej wydzielono 10\,\% obserwacji na walidację.
Ze względu na naturę szeregu czasowego podczas uczenia nie stosowano losowego mieszania próbek (\texttt{shuffle=False}).

Model trenowano algorytmem Adam z początkową szybkością uczenia $5 \times 10^{-5}$,
funkcją straty \texttt{binary\_crossentropy} oraz wagami klas (\texttt{compute\_class\_weight})
kompensującymi niezbalansowanie etykiet.
Zastosowano wczesne zatrzymanie (\emph{early stopping}) monitorujące dokładność walidacyjną
oraz redukcję szybkości uczenia przy stagnacji straty.
Predykcje ciągłe z warstwy sigmoid konwertowano na klasy za pomocą progu $0{,}5$.

Jakość modelu oceniano metrykami: accuracy, precision, recall, F1 oraz macierzą pomyłek.
Wyniki porównywano z naiwnymi strategiami bazowymi przewidującymi zawsze wzrost lub zawsze spadek kursu.

\section*{Zakres i ograniczenia}

Praca koncentruje się wyłącznie na prognozie \emph{kierunku} zmiany ceny, a nie na estymacji jej poziomu
ani wielkości zwrotu. Nie uwzględniono kosztów transakcyjnych, spreadu bid-ask ani dźwigni finansowej,
co ogranicza bezpośrednią aplikowalność wyników w realnym handlu algorytmicznym.

Model wytrenowano na danych historycznych z okresu od 2020~roku; jego skuteczność w przyszłości
nie jest gwarantowana ze względu na ryzyko przeuczenia oraz możliwość zmiany reżimu rynkowego
(np.\ interwencje banków centralnych, kryzysy geopolityczne). Wyniki należy traktować jako eksperyment badawczy,
a nie rekomendację inwestycyjną.

\section*{Struktura pracy}

% --- Dostosuj nazwy rozdziałów do własnego spisu treści ---

Praca składa się z następujących rozdziałów:
\begin{enumerate}
    \item Wstęp --- motywacja, problem badawczy, cel i zakres pracy.
    \item Przegląd literatury i podstawy teoretyczne --- szeregi czasowe finansowe,
          sieci LSTM, architektura Transformer oraz wybrane wskaźniki analizy technicznej.
    \item Przygotowanie danych i inżynieria cech --- opis źródeł danych, preprocessingu
          i konstrukcji 24~zmiennych wejściowych.
    \item Architektura i trenowanie modelu --- szczegółowy opis hybrydy LSTM\,+\,Transformer,
          hiperparametrów i procedury walidacji.
    \item Wyniki i analiza --- metryki klasyfikacji, macierz pomyłek, porównanie z benchmarkami
          oraz wizualizacja sygnałów kupna i sprzedaży.
    \item Podsumowanie i wnioski --- ocena osiągniętych rezultatów i kierunki dalszych badań.
\end{enumerate}

% =============================================================================

% --- Odkomentuj odpowiednią linię lub umieść treść poniżej pod własnym nagłówkiem ---

% \chapter{Wstęp}
% \section{Wstęp}

Rynek walutowy stanowi jeden z najbardziej płynnych i globalnie zintegrowanych segmentów rynków finansowych.
Kurs wymiany dolara amerykańskiego względem złotego polskiego (USD/PLN) ma bezpośrednie znaczenie dla polskiej gospodarki:
wpływa na koszty importu, wartość zobowiązań zadłużonych w walucie obcej oraz poziom inflacji importowanej.
Instytucje finansowe, przedsiębiorstwa i inwestorzy indywidualni poszukują narzędzi wspierających podejmowanie decyzji na tym rynku,
jednak prognozowanie kursów walutowych pozostaje zadaniem wyjątkowo trudnym.

Szeregi czasowe notowań finansowych charakteryzują się wysokim poziomem szumu, nieliniowością zależności
oraz okresową zmianą reżimów rynkowych. Klasyczne modele ekonometryczne, oparte na założeniach stacjonarności
i liniowości, często nie wystarczają do uchwycenia złożonych wzorców zachowania cen.
W ostatnich latach coraz większą popularnością cieszą się metody uczenia głębokiego, w szczególności sieci rekurencyjne
typu LSTM (Long Short-Term Memory) oraz architektury oparte na mechanizmie uwagi (attention), takie jak Transformer.
Pozwalają one na modelowanie wielowymiarowych sekwencji czasowych i łączenie informacji z różnych horyzontów temporalnych.

\section*{Problem badawczy}

Niniejsza praca podejmuje problem prognozowania kierunku zmiany kursu USD/PLN w horyzoncie 14~dni roboczych.
Główne pytanie badawcze brzmi: \emph{czy na podstawie 30-dniowej historii cech technicznych
i makroekonomicznych można skutecznie przewidzieć, czy kurs wzrośnie, czy spadnie w ciągu kolejnych 14~dni?}

Punktem wyjścia jest skumulowany zwrot procentowy kursu zamknięcia w horyzoncie prognozy.
Ze względu na praktyczną interpretację wyników (sygnały kupna i sprzedaży) zadanie sprowadzono
do klasyfikacji binarnej: klasa~1 oznacza wzrost kursu, klasa~0 oznacza spadek.
Etykieta dla dnia $t$ definiowana jest jako:
\begin{equation}
y_t =
\begin{cases}
1 & \text{gdy } \dfrac{P_{t+14}}{P_t} - 1 > 0, \\[6pt]
0 & \text{gdy } \dfrac{P_{t+14}}{P_t} - 1 \le 0.
\end{cases}
\label{eq:target}
\end{equation}
gdzie $P_t$ oznacza kurs zamknięcia pary USD/PLN w dniu~$t$.

Skumulowany zwrot w procentach wyraża się natomiast wzorem:
\begin{equation}
    r_t = \left(\frac{P_{t+14}}{P_t} - 1\right) \cdot 100\,\%.
    \label{eq:return}
\end{equation}

\section*{Cel pracy}

Celem pracy jest zaprojektowanie, implementacja i ocena hybrydowego modelu głębokiego łączącego warstwę LSTM
z blokiem Transformer w zadaniu klasyfikacji kierunku kursu USD/PLN.
W ramach realizacji celu:
\begin{itemize}
    \item przygotowano wielowymiarowy zestaw cech obejmujący wskaźniki analizy technicznej
          oraz stopy procentowe banków centralnych (Fed, EBC, NBP),
    \item zbudowano i wytrenowano model na historycznych danych dziennych,
    \item oceniono skuteczność predykcji za pomocą metryk klasyfikacji
          (accuracy, precision, recall, F1) oraz porównano wyniki z prostymi benchmarkami
          (strategia ,,zawsze wzrost'' i ,,zawsze spadek'').
\end{itemize}

\section*{Opis metody}

\subsection*{Dane wejściowe}

Analizie poddano dzienne notowania z pliku \texttt{dane\_ekonomiczne.csv}, którego źródłem jest Yahoo Finance oraz World Bank Data.
Dane obejmują m.in.\ ceny OHLC pary USD/PLN, skorelowane pary walutowe (EUR/PLN, EUR/USD)
oraz stopy procentowe banków centralnych Stanów Zjednoczonych, strefy euro i Polski.
Z zestawu usunięto obserwacje przypadające na weekendy, aby zachować ciągłość danych sesyjnych.
Przed przekazaniem do modelu wszystkie cechy numeryczne standaryzowano za pomocą \texttt{StandardScaler},
dopasowanego wyłącznie na zbiorze treningowym.

\subsection*{Inżynieria cech}

Dla każdego dnia obliczono łącznie 24~zmienne wejściowe, w tym:
\begin{itemize}
    \item dzienne stopy zwrotu kursów zamknięcia (USD/PLN, EUR/PLN, EUR/USD),
    \item wskaźniki momentum za okresy 3, 5 i 10~dni,
    \item RSI (14~okresów), ATR (14~okresów), pozycję i szerokość pasm Bollingera (okno 20~dni),
    \item wskaźniki MACD wraz z linią sygnałową i histogramem,
    \item zmienność kroczącą, $z$-score (okno 20~dni), dzienny zakres High/Low oraz stosunek High/Low,
    \item stopy procentowe: US\_FED, EA\_ECB, PL\_NBP,
    \item opóźnione kierunki zmian ceny (lagi 1, 2, 3 i 5~dni).
\end{itemize}

\subsection*{Architektura modelu}

Model przyjmuje na wejściu sekwencję $X_t \in \mathbb{R}^{30 \times 24}$,
tj.\ 30~kolejnych dni obserwacji po 24~cechy (okno przesuwne).
Architektura hybrydowa składa się z następujących etapów:
\begin{enumerate}
    \item warstwa LSTM z 128~jednostkami ukrytymi (\texttt{return\_sequences=True}),
          wychwytująca lokalne zależności czasowe w sekwencji,
    \item warstwa \texttt{Dropout} ($p = 0{,}3$) redukująca przeuczenie,
    \item kodowanie pozycyjne (\emph{Positional Encoding}) nadające informację o położeniu obserwacji w oknie,
    \item blok Transformer z 8~głowicami uwagi (\emph{Multi-Head Attention}),
          $d_{\text{model}} = 128$ i warstwą sprzężeń zwrotnych o wymiarze 256,
          modelujący długoterminowe zależności w obrębie 30-dniowego okna,
    \item agregacja \texttt{GlobalAveragePooling1D}, po której następują warstwy gęste
          (64 i 16~neuronów, aktywacja ReLU, regularyzacja L2) oraz warstwa wyjściowa
          z aktywacją sigmoid (jedna neurona, klasyfikacja binarna).
\end{enumerate}

Połączenie LSTM i Transformera uzasadnia fakt, że sieć rekurencyjna efektywnie koduje krótkoterminowe wzorce cenowe,
natomiast mechanizm self-attention pozwala ważyć znaczenie poszczególnych dni w historycznym oknie,
co jest istotne przy analizie wielotygodniowych trendów.

\subsection*{Trenowanie i ewaluacja}

Dane podzielono chronologicznie: 80\,\% stanowi zbiór treningowy, 20\,\%~--- testowy.
Z części treningowej wydzielono 10\,\% obserwacji na walidację.
Ze względu na naturę szeregu czasowego podczas uczenia nie stosowano losowego mieszania próbek (\texttt{shuffle=False}).

Model trenowano algorytmem Adam z początkową szybkością uczenia $5 \times 10^{-5}$,
funkcją straty \texttt{binary\_crossentropy} oraz wagami klas (\texttt{compute\_class\_weight})
kompensującymi niezbalansowanie etykiet.
Zastosowano wczesne zatrzymanie (\emph{early stopping}) monitorujące dokładność walidacyjną
oraz redukcję szybkości uczenia przy stagnacji straty.
Predykcje ciągłe z warstwy sigmoid konwertowano na klasy za pomocą progu $0{,}5$.

Jakość modelu oceniano metrykami: accuracy, precision, recall, F1 oraz macierzą pomyłek.
Wyniki porównywano z naiwnymi strategiami bazowymi przewidującymi zawsze wzrost lub zawsze spadek kursu.

\section*{Zakres i ograniczenia}

Praca koncentruje się wyłącznie na prognozie \emph{kierunku} zmiany ceny, a nie na estymacji jej poziomu
ani wielkości zwrotu. Nie uwzględniono kosztów transakcyjnych, spreadu bid-ask ani dźwigni finansowej,
co ogranicza bezpośrednią aplikowalność wyników w realnym handlu algorytmicznym.

Model wytrenowano na danych historycznych z okresu od 2020~roku; jego skuteczność w przyszłości
nie jest gwarantowana ze względu na ryzyko przeuczenia oraz możliwość zmiany reżimu rynkowego
(np.\ interwencje banków centralnych, kryzysy geopolityczne). Wyniki należy traktować jako eksperyment badawczy,
a nie rekomendację inwestycyjną.

\section*{Struktura pracy}

% --- Dostosuj nazwy rozdziałów do własnego spisu treści ---

Praca składa się z następujących rozdziałów:
\begin{enumerate}
    \item Wstęp --- motywacja, problem badawczy, cel i zakres pracy.
    \item Przegląd literatury i podstawy teoretyczne --- szeregi czasowe finansowe,
          sieci LSTM, architektura Transformer oraz wybrane wskaźniki analizy technicznej.
    \item Przygotowanie danych i inżynieria cech --- opis źródeł danych, preprocessingu
          i konstrukcji 24~zmiennych wejściowych.
    \item Architektura i trenowanie modelu --- szczegółowy opis hybrydy LSTM\,+\,Transformer,
          hiperparametrów i procedury walidacji.
    \item Wyniki i analiza --- metryki klasyfikacji, macierz pomyłek, porównanie z benchmarkami
          oraz wizualizacja sygnałów kupna i sprzedaży.
    \item Podsumowanie i wnioski --- ocena osiągniętych rezultatów i kierunki dalszych badań.
\end{enumerate}




% ============================================================================

\chapter{Wstęp teoretyczny}

% ============================================================================



\section{Szeregi czasowe finansowe}



Notowania instrumentów finansowych tworzą szereg czasowy --- uporządkowany ciąg obserwacji

$(x_1, x_2, \ldots, x_T)$, w którym kolejność ma znaczenie predykcyjne.

Kurs walutowy USD/PLN jest procesem stochastycznym o wysokiej zmienności,

często opisywanym jako ,,losowy spacer'' z dryftem (hipoteza efektywnego rynku w słabej formie).

Mimo tego w danych historycznych można zaobserwować krótkoterminowe zależności autokorelacyjne

oraz klastry zmienności, które stanowią podstawę analizy technicznej i modeli uczących się.



\section{Sieci LSTM}



LSTM (Long Short-Term Memory) to architektura rekurencyjnej sieci neuronowej (RNN),

zaproponowana przez Hochreitera i Schmidhubera w~1997~roku.

W przeciwieństwie do klasycznego RNN, LSTM rozwiązuje problem zanikającego gradientu

dzięki komórce pamięci i bramkom: wejścia (\emph{input}), zapomnienia (\emph{forget})

oraz wyjścia (\emph{output}).



Dla sekwencji wejściowej $\mathbf{x}_t$ stan ukryty $\mathbf{h}_t$ aktualizowany jest rekurencyjnie:

\begin{equation}

    \mathbf{f}_t = \sigma(\mathbf{W}_f [\mathbf{h}_{t-1}, \mathbf{x}_t] + \mathbf{b}_f), \quad

    \mathbf{i}_t = \sigma(\mathbf{W}_i [\mathbf{h}_{t-1}, \mathbf{x}_t] + \mathbf{b}_i),

\end{equation}

\begin{equation}

    \tilde{\mathbf{c}}_t = \tanh(\mathbf{W}_c [\mathbf{h}_{t-1}, \mathbf{x}_t] + \mathbf{b}_c), \quad

    \mathbf{c}_t = \mathbf{f}_t \odot \mathbf{c}_{t-1} + \mathbf{i}_t \odot \tilde{\mathbf{c}}_t,

\end{equation}

\begin{equation}

    \mathbf{o}_t = \sigma(\mathbf{W}_o [\mathbf{h}_{t-1}, \mathbf{x}_t] + \mathbf{b}_o), \quad

    \mathbf{h}_t = \mathbf{o}_t \odot \tanh(\mathbf{c}_t),

\end{equation}

gdzie $\sigma$ oznacza funkcję sigmoid, a $\odot$ iloczyn Hadamarda.

W projekcie warstwa LSTM z 128~jednostkami ukrytymi koduje lokalne wzorce cenowe

z 30-dniowego okna obserwacji.



\section{Architektura Transformer i mechanizm uwagi}



Transformer, wprowadzony przez Vaswaniego \emph{et al.} (2017), opiera się na mechanizmie

\emph{self-attention}, który pozwala każdemu elementowi sekwencji ważyć znaczenie pozostałych.

Dla macierzy zapytań $\mathbf{Q}$, kluczy $\mathbf{K}$ i wartości $\mathbf{V}$:

\begin{equation}

    \mathrm{Attention}(\mathbf{Q}, \mathbf{K}, \mathbf{V}) =

    \mathrm{softmax}\!\left(\frac{\mathbf{Q}\mathbf{K}^\top}{\sqrt{d_k}}\right)\mathbf{V}.

\end{equation}



W implementacji (\texttt{TransformerBlock}) zastosowano wielogłowicową uwagę

(\emph{Multi-Head Attention}, 8~głów) oraz kodowanie pozycyjne sinusoidalne

(\texttt{PositionalEncoding}), które nadaje informację o położeniu dnia w oknie czasowym.

Połączenie LSTM\,+\,Transformer pozwala jednocześnie modelować krótkoterminową dynamikę

oraz długoterminowe zależności w obrębie 30~dni.



\section{Wybrane wskaźniki analizy technicznej}



Tabela~\ref{tab:wskazniki} zestawia wskaźniki użyte jako cechy modelu.



\begin{table}[H]

\centering

\caption{Wskaźniki techniczne wykorzystane w \texttt{LSTM2.py}}

\label{tab:wskazniki}

\begin{tabular}{@{}llp{7cm}@{}}

\toprule

\textbf{Wskaźnik} & \textbf{Okno} & \textbf{Opis} \\

\midrule

Momentum & 3, 5, 10 dni & Procentowa zmiana ceny zamknięcia \\

RSI & 14 & Siła relatywna: $100 - 100/(1 + RS)$ \\

ATR & 14 & Średni prawdziwy zakres (volatility) \\

Bollinger & 20 & Pozycja ceny względem pasm; szerokość pasm \\

MACD & 12/26/9 & Różnica EMA, linia sygnałowa, histogram \\

$z$-score & 20 & Standaryzacja ceny względem średniej kroczącej \\

Zmienność & 10 & Odchylenie std.\ dziennych stóp zwrotu \\

\bottomrule

\end{tabular}

\end{table}



\section{Zadanie klasyfikacji binarnej}



Zamiast prognozować dokładną wartość kursu, model przewiduje \emph{kierunek} zmiany

w horyzoncie $H = 14$~dni roboczych. Etykieta definiowana jest jako:

\begin{equation}

    y_t = \mathbb{1}\!\left[\frac{P_{t+H}}{P_t} - 1 > 0\right],

    \label{eq:target-doc}

\end{equation}

gdzie $P_t$ to kurs zamknięcia USD/PLN. Klasa~1 oznacza wzrost (sygnał LONG/KUP),

klasa~0 --- spadek (sygnał SHORT/SPRZEDAJ).



% ============================================================================

\chapter{Opis działania aplikacji}

% ============================================================================



\section{Uwagi przed uruchomieniem}



\begin{itemize}

    \item Skrypt \texttt{LSTM2.py} wymaga pliku \texttt{dane\_ekonomiczne.csv}

          w tym samym katalogu co skrypt.

    \item W pliku CSV kolumny stóp procentowych noszą nazwy

          \texttt{US\_FED\_IR}, \texttt{EA\_ECB\_IR}, \texttt{PL\_NBP\_IR},

          natomiast w kodzie filtrowane są jako \texttt{US\_FED}, \texttt{EA\_ECB}, \texttt{PL\_NBP}.

          Przed uruchomieniem należy dodać rename kolumn lub ujednolicić nazwy w kodzie.

    \item Zalecane środowisko: Python~3.10+, TensorFlow~2.x, pandas, scikit-learn, matplotlib.

    \item Ze względu na losowość inicjalizacji wag sieci ustawiono stałe ziarno (\texttt{SEED = 42})

          dla powtarzalności wyników.

    \item Trenowanie na CPU może trwać kilkanaście--kilkadziesiąt minut; GPU znacząco przyspiesza proces.

    \item Wykresy (\texttt{plt.show()}) wymagają środowiska graficznego lub backendu

          \texttt{Agg} z zapisem do pliku.

\end{itemize}



\section{O aplikacji}



Aplikacja została napisana w języku Python z wykorzystaniem biblioteki TensorFlow/Keras.

Główny plik implementacyjny to \texttt{LSTM2.py}. Powiązane pliki w repozytorium:

\begin{itemize}

    \item \texttt{DataDownloader.py} --- pobieranie danych z Yahoo Finance i World Bank,

    \item \texttt{DataTransformer.py} --- transformacja danych do formatu dziennego,

    \item \texttt{LSTM.py} --- wcześniejsza, prostsza wersja modelu (regresja),

    \item \texttt{lstm\_test.py} --- wariant eksperymentalny z innymi hiperparametrami.

\end{itemize}



\section{Przepływ danych}



Schemat przetwarzania danych w \texttt{LSTM2.py}:



\begin{enumerate}

    \item Wczytanie \texttt{dane\_ekonomiczne.csv}.

    \item Usunięcie zbędnych kolumn i weekendów.

    \item Obliczenie 24~cech technicznych i makroekonomicznych.

    \item Usunięcie wierszy z wartościami \texttt{NaN} (okna kroczące).

    \item Podział chronologiczny 80\,\% / 20\,\% (trening / test).

    \item Standaryzacja \texttt{StandardScaler} (dopasowanie tylko na treningu).

    \item Budowa okien przesuwnych $30 \times 24$.

    \item Trenowanie modelu, predykcja, ewaluacja metryk.

\end{enumerate}



\section{Parametry konfiguracyjne}



Tabela~\ref{tab:parametry} zestawia główne hiperparametry modelu.



\begin{table}[H]

\centering

\caption{Hiperparametry modelu w \texttt{LSTM2.py}}

\label{tab:parametry}

\begin{tabular}{@{}lll@{}}

\toprule

\textbf{Parametr} & \textbf{Wartość} & \textbf{Opis} \\

\midrule

\texttt{sliding\_window} & 30 & Długość okna wejściowego (dni) \\

\texttt{prediction\_horizon} & 14 & Horyzont prognozy (dni) \\

\texttt{num\_of\_epochs} & 50 & Maksymalna liczba epok \\

LSTM units & 128 & Jednostki ukryte warstwy LSTM \\

\texttt{num\_heads} & 8 & Liczba głów uwagi w Transformer \\

\texttt{embed\_dim} & 128 & Wymiar osadzenia \\

\texttt{ff\_dim} & 256 & Wymiar warstwy FFN \\

Dropout & 0.3 / 0.5 / 0.3 & Po LSTM / Dense64 / Dense16 \\

Learning rate & $5 \times 10^{-5}$ & Szybkość uczenia (Adam) \\

Batch size & 16 & Rozmiar minibatcha \\

Podział train/test & 80\,\% / 20\,\% & Chronologiczny \\

Podział train/val & 90\,\% / 10\,\% & Z części treningowej \\

Próg klasyfikacji & 0.5 & Konwersja sigmoid $\rightarrow$ klasa \\

\bottomrule

\end{tabular}

\end{table}



\section{Zestaw cech wejściowych}



Model wykorzystuje 24~zmienne (kształt wejścia: $30 \times 24$):



\begin{enumerate}

    \item \textbf{Stopy zwrotu:} \texttt{Close\_USDPLN}, \texttt{Close\_EURPLN}, \texttt{Close\_EURUSD}

    \item \textbf{Momentum:} \texttt{mom\_3}, \texttt{mom\_5}, \texttt{mom\_10}

    \item \textbf{Techniczne:} \texttt{RSI}, \texttt{ATR}, \texttt{bb\_pos}, \texttt{bb\_width},

          \texttt{usd\_volatility}, \texttt{zscore\_20}, \texttt{MACD}, \texttt{MACD\_signal},

          \texttt{MACD\_histogram}

    \item \textbf{High/Low:} \texttt{usd\_daily\_range}, \texttt{usd\_hl\_ratio}

    \item \textbf{Stopy procentowe:} \texttt{US\_FED}, \texttt{EA\_ECB}, \texttt{PL\_NBP}

    \item \textbf{Opóźniony kierunek:} \texttt{dir\_lag\_1}, \texttt{dir\_lag\_2},

          \texttt{dir\_lag\_3}, \texttt{dir\_lag\_5}

\end{enumerate}



\section{Architektura modelu}



Model hybrydowy składa się z następujących warstw (kolejność w \texttt{LSTM2.py}):



\begin{lstlisting}

Input(shape=(30, 24))

  -> LSTM(128, return_sequences=True)

  -> Dropout(0.3)

  -> PositionalEncoding(max_len=30, embed_dim=128)

  -> TransformerBlock(embed_dim=128, num_heads=8, ff_dim=256)

  -> GlobalAveragePooling1D()

  -> Dense(64, relu, L2=0.005) -> Dropout(0.5)

  -> Dense(16, relu, L2=0.005) -> Dropout(0.3)

  -> Dense(1, sigmoid)   # wyjscie binarne

\end{lstlisting}



Funkcja straty: \texttt{binary\_crossentropy}.

Optymalizator: Adam z \texttt{ReduceLROnPlateau} i \texttt{EarlyStopping}

(monitorowanie \texttt{val\_accuracy}, \texttt{patience=50}).



\section{Uruchomienie aplikacji}



\begin{enumerate}

    \item Przejdź do katalogu \texttt{Kod/}:

\begin{lstlisting}

cd "c:\Users\jakja\Desktop\Studia\4Sem\Sztuczna inteligencja projekt\Kod"

\end{lstlisting}

    \item Zainstaluj zależności (jeśli brakuje):

\begin{lstlisting}

pip install tensorflow pandas numpy scikit-learn matplotlib

\end{lstlisting}

    \item Uruchom skrypt:

\begin{lstlisting}

python LSTM2.py

\end{lstlisting}

    \item Po zakończeniu trenowania w konsoli pojawią się:

          bilans klas, rozmiary zbiorów, \texttt{model.summary()},

          tabela wyników, macierz pomyłek, metryki klasyfikacji,

          porównanie z benchmarkami oraz wykres sygnałów KUP/SPRZEDAJ.

\end{enumerate}



\section{Interpretacja wyników na wyjściu konsoli}



Program wypisuje następujące sekcje (funkcje pomocnicze na końcu pliku):



\begin{itemize}

    \item \textbf{Class balance} --- procent klas UP/DOWN w zbiorze treningowym.

    \item \textbf{Class weights} --- wagi kompensujące niezbalansowanie.

    \item \textbf{Tabela wyników} --- data, cena, rzeczywisty i przewidziany kierunek.

    \item \textbf{Macierz błędu} --- TP, TN, FP, FN.

    \item \textbf{Metryki} --- accuracy, precision, recall, F1.

    \item \textbf{Benchmark} --- porównanie z strategiami ,,zawsze LONG'' i ,,zawsze SHORT''.

    \item \textbf{Wykres} --- notowania USD/PLN z oznaczeniem sygnałów (zielone $\wedge$ = KUP,

          czerwone $\vee$ = SPRZEDAJ).

\end{itemize}



% ============================================================================

\chapter{Badania i testy modelu}

% ============================================================================



\section{Cel rozdziału}



Niniejszy rozdział opisuje, \emph{gdzie}, \emph{jak} i \emph{co} testować,

wzorując się na strukturze badań eksperymentalnych z dokumentacji projektowej

(rozdział ,,Badania'' w pracy referencyjnej PSO).

Celem testów jest wykazanie, że model:

\begin{enumerate}

    \item uczy się sensownych wzorców (a nie tylko przeucza się na danych treningowych),

    \item przewyższa proste strategie bazowe,

    \item zachowuje stabilność przy zmianie wybranych hiperparametrów.

\end{enumerate}



\section{Gdzie przeprowadzać testy}



\subsection{Testy wbudowane w \texttt{LSTM2.py} (obowiązkowe)}



Główny skrypt zawiera już podstawową ewaluację na zbiorze testowym (ostatnie 20\,\% danych

chronologicznie). Po uruchomieniu \texttt{python LSTM2.py} otrzymujesz:

\begin{itemize}

    \item metryki na \textbf{zbiorze testowym} (out-of-sample),

    \item benchmark naiwny,

    \item wizualizację sygnałów.

\end{itemize}

To jest \textbf{główny test końcowy} --- wyniki z tego zbioru należy opisać w dokumentacji i wnioskach.



\subsection{Walidacja w trakcie trenowania}



Podzbiór walidacyjny (10\,\% z danych treningowych) służy do:

\begin{itemize}

    \item monitorowania \texttt{val\_accuracy} i \texttt{val\_loss},

    \item wczesnego zatrzymania (\texttt{EarlyStopping}),

    \item redukcji learning rate (\texttt{ReduceLROnPlateau}).

\end{itemize}

\textbf{Nie} należy traktować metryk walidacyjnych jako ostatecznej oceny modelu ---

służą one do strojenia i zatrzymania uczenia.



\subsection{Testy dodatkowe (zalecane w pracy)}



\begin{table}[H]

\centering

\caption{Mapa testów --- gdzie i co sprawdzać}

\label{tab:testy}

\begin{tabular}{@{}p{3.5cm}p{4cm}p{6cm}@{}}

\toprule

\textbf{Miejsce} & \textbf{Plik / metoda} & \textbf{Co testujesz} \\

\midrule

Zbiór testowy & \texttt{LSTM2.py} (linie 333--463) & Jakość predykcji out-of-sample \\

Walidacja & \texttt{model.fit(..., validation\_data=...)} & Przeuczenie, early stopping \\

Zbiór treningowy & \texttt{training\_predictions\_lstm} & Overfitting (porównaj z testem) \\

Eksperymenty & modyfikacja parametrów w \texttt{LSTM2.py} & Wrażliwość na hiperparametry \\

Walk-forward & osobny skrypt / \texttt{lstm\_test.py} & Stabilność w czasie \\

\bottomrule

\end{tabular}

\end{table}



\section{Jak podejść do testów --- metodologia}



\subsection{Zasada podziału chronologicznego}



Dane finansowe mają strukturę czasową --- \textbf{nie wolno} losować próbek

(\texttt{shuffle=False} jest poprawne).

Podział 80/20 chronologiczny symuluje sytuację, w której model trenuje na przeszłości

i prognozuje przyszłość. Jest to minimalny wymóg poprawności metodologicznej.



\subsection{Metryki do raportowania}



Dla klasyfikacji binarnej kierunku należy podać:



\begin{table}[H]

\centering

\caption{Metryki ewaluacji}

\label{tab:metryki}

\begin{tabular}{@{}ll@{}}

\toprule

\textbf{Metryka} & \textbf{Interpretacja} \\

\midrule

Accuracy & Odsetek trafnych predykcji kierunku \\

Precision (klasa 1) & Jak często sygnał KUP był trafny \\

Recall (klasa 1) & Jaki odsetek wzrostów model wykrył \\

F1 & Średnia harmoniczna precision i recall \\

Macierz pomyłek & TP, TN, FP, FN --- szczegółowa analiza błędów \\

\bottomrule

\end{tabular}

\end{table}



\subsection{Benchmarki obowiązkowe}



Funkcja \texttt{naive\_benchmark()} porównuje model ze strategiami:

\begin{itemize}

    \item \textbf{Zawsze LONG} --- zawsze przewiduj wzrost,

    \item \textbf{Zawsze SHORT} --- zawsze przewiduj spadek.

\end{itemize}

Jeśli accuracy modelu nie przewyższa ,,zawsze LONG'', model nie wnosi wartości predykcyjnej

(dla niezbalansowanych danych z przewagą wzrostów ta strategia może mieć wysoką accuracy).



\subsection{Test przeuczenia (overfitting)}



Porównaj metryki:

\begin{itemize}

    \item \textbf{Trening:} \texttt{accuracy\_score(Y\_train\_binary, training\_predictions\_lstm)}

    \item \textbf{Test:} \texttt{accuracy\_score(Y\_test\_binary, predictions\_lstm)}

\end{itemize}

Duża różnica (np.\ trening 90\,\%, test 52\,\%) wskazuje na przeuczenie.

W pracy opisz tę różnicę i zaproponuj środki zaradcze (większy dropout, mniej epok, regularyzacja L2).



\section{Co uwzględnić w badaniach --- scenariusze testowe}



Poniższe scenariusze odpowiadają rozdziałowi ,,Badania'' z dokumentacji referencyjnej PSO

--- zamiast testować różne funkcje celu, testujesz różne konfiguracje modelu.



\subsection{Scenariusz 1: Konfiguracja bazowa (domyślna)}



\begin{itemize}

    \item \texttt{sliding\_window = 30}, \texttt{prediction\_horizon = 14}

    \item Architektura LSTM(128) + Transformer(8 heads)

    \item \texttt{learning\_rate = 5e-5}, \texttt{batch\_size = 16}

\end{itemize}

\textbf{Cel:} uzyskanie wyniku referencyjnego do porównań.

Zapisz: accuracy, precision, recall, F1, macierz pomyłek, wykres sygnałów.



\subsection{Scenariusz 2: Wpływ długości okna czasowego}



Zmień \texttt{sliding\_window} na wartości: 15, 30, 60 (trzy osobne uruchomienia).



\begin{table}[H]

\centering

\caption{Szablon tabeli wyników --- wpływ okna czasowego}

\label{tab:okno}

\begin{tabular}{@{}cccc@{}}

\toprule

\textbf{Okno (dni)} & \textbf{Accuracy} & \textbf{F1} & \textbf{Uwagi} \\

\midrule

15  & --- & --- & Krótsza pamięć \\

30  & --- & --- & Konfiguracja bazowa \\

60  & --- & --- & Dłuższa pamięć \\

\bottomrule

\end{tabular}

\end{table}



\textbf{Hipoteza:} zbyt krótkie okno nie uchwyci trendu; zbyt długie wprowadzi szum.



\subsection{Scenariusz 3: Wpływ horyzontu prognozy}



Zmień \texttt{prediction\_horizon} na: 5, 14, 21~dni.



\textbf{Cel:} sprawdzenie, czy model lepiej przewiduje krótki czy średni horyzont.

Krótszy horyzont zwykle daje wyższą accuracy, ale ma mniejszą użyteczność inwestycyjną.



\subsection{Scenariusz 4: Wpływ rozmiaru warstwy LSTM}



Porównaj \texttt{LSTM(64)} vs \texttt{LSTM(128)} vs \texttt{LSTM(256)}.



\textbf{Cel:} analiza kompromisu między zdolnością uczenia a ryzykiem przeuczenia.



\subsection{Scenariusz 5: Model bez Transformera (ablation study)}



Tymczasowo usuń lub zakomentuj warstwy \texttt{PositionalEncoding} i \texttt{TransformerBlock},

zastępując je np.\ \texttt{GlobalAveragePooling1D} bezpośrednio po LSTM.



\textbf{Cel:} wykazanie, czy hybryda LSTM\,+\,Transformer daje lepsze wyniki niż samo LSTM.

Jest to analogia do porównania funkcji Ackley'a i Wielbłąda w dokumentacji PSO.



\subsection{Scenariusz 6: Zakres dat (opcjonalnie)}



Odkomentuj filtr dat w \texttt{LSTM2.py} (linie 44--45), np.:

\begin{lstlisting}

data = data[(data['Date'] >= "2022-01-01") &

            (data['Date'] <= "2023-12-31")]

\end{lstlisting}

Porównaj wyniki dla różnych reżimów rynkowych (pandemia, inflacja, podwyżki stóp).



\section{Procedura przeprowadzenia pojedynczego testu}



\begin{enumerate}

    \item Ustaw \texttt{SEED = 42} (powtarzalność).

    \item Zmień \textbf{jeden} hiperparametr na raz (kontrolowany eksperyment).

    \item Uruchom \texttt{python LSTM2.py} i poczekaj na zakończenie trenowania.

    \item Zapisz wyniki do tabeli (Excel/CSV/LaTeX) --- accuracy, precision, recall, F1,

          TP/TN/FP/FN, czas trenowania.

    \item Zapisz wykres sygnałów jako plik PNG (dodaj \texttt{plt.savefig()} przed \texttt{show()}).

    \item Porównaj z konfiguracją bazową i benchmarkiem naiwnym.

    \item Sformułuj wniosek dla danego scenariusza (1--3 zdania).

\end{enumerate}



\section{Szablon tabeli zbiorczej wyników}



\begin{table}[H]

\centering

\caption{Przykładowa tabela zbiorcza wyników badań (do uzupełnienia)}

\label{tab:zbiorcza}

\begin{tabular}{@{}lccccc@{}}

\toprule

\textbf{Scenariusz} & \textbf{Acc.} & \textbf{Prec.} & \textbf{Rec.} & \textbf{F1} & \textbf{vs LONG} \\

\midrule

Bazowy (okno=30, H=14)       & --- & --- & --- & --- & --- \\

Okno = 15                    & --- & --- & --- & --- & --- \\

Okno = 60                    & --- & --- & --- & --- & --- \\

Horyzont = 5                 & --- & --- & --- & --- & --- \\

Horyzont = 21                & --- & --- & --- & --- & --- \\

LSTM(64) bez Transformer     & --- & --- & --- & --- & --- \\

\bottomrule

\end{tabular}

\end{table}



\section{Na co uważać --- typowe błędy}



\begin{itemize}

    \item \textbf{Data leakage} --- skalowanie na całym zbiorze zamiast tylko na treningu

          (w kodzie jest poprawnie: \texttt{fit\_transform} tylko na \texttt{train\_data}).

    \item \textbf{Losowe mieszanie} --- \texttt{shuffle=True} zaburzyłoby kolejność czasową.

    \item \textbf{Wielokrotne testowanie na zbiorze testowym} --- po wielu eksperymentach

          dostosowanych do testu, wynik przestaje być uczciwy; trzymaj zbiór testowy

          ,,na koniec'' lub wydziel trzeci zbiór hold-out.

    \item \textbf{Interpretacja samej accuracy} --- przy niezbalansowanych klasach

          accuracy może być myląca; zawsze podawaj macierz pomyłek i F1.

    \item \textbf{Brak benchmarku} --- bez porównania z ,,zawsze LONG'' nie można ocenić,

          czy model w ogóle się nauczył.

\end{itemize}



% ============================================================================

\chapter{Wnioski}

% ============================================================================



% --- Uzupełnij po przeprowadzeniu badań ---



Po przeprowadzeniu badań opisanych w rozdziale~\ref{tab:testy} należy sformułować wnioski.

Poniżej szablon punktów do uzupełnienia własnymi wynikami:



\begin{enumerate}

    \item \textbf{Skuteczność modelu bazowego} --- czy hybryda LSTM\,+\,Transformer

          przewyższa benchmark ,,zawsze LONG'' na zbiorze testowym?

          Jaka jest wartość accuracy, precision, recall i F1?

    \item \textbf{Wpływ hiperparametrów} --- która długość okna i horyzont prognozy

          dawały najlepsze wyniki? Czy różnice są istotne?

    \item \textbf{Rola Transformera} --- czy ablation study (scenariusz~5) wykazał

          poprawę względem samego LSTM?

    \item \textbf{Przeuczenie} --- jaka jest różnica metryk między zbiorem treningowym

          a testowym? Czy regularyzacja (dropout, L2, early stopping) była wystarczająca?

    \item \textbf{Ograniczenia} --- model przewiduje tylko kierunek (nie wielkość zwrotu),

          nie uwzględnia kosztów transakcji; wyniki historyczne nie gwarantują

          skuteczności w przyszłości.

    \item \textbf{Kierunki rozwoju} --- np.\ walk-forward validation, więcej par walutowych,

          integracja z modułem GUI (\texttt{main\_gui.py}), prognoza probabilistyczna.

\end{enumerate}



% ============================================================================

\chapter{Bibliografia}

% ============================================================================



\begin{enumerate}

    \item Hochreiter S., Schmidhuber J.,

          \emph{Long Short-Term Memory},

          Neural Computation, 9(8), 1735--1780, 1997.

    \item Vaswani A., Shazeer N., Parmar N., \emph{et al.},

          \emph{Attention Is All You Need},

          Advances in Neural Information Processing Systems, 2017.

    \item Murphy J. J.,

          \emph{Technical Analysis of the Financial Markets},

          New York Institute of Finance, 1999.

    \item Goodfellow I., Bengio Y., Courville A.,

          \emph{Deep Learning},

          MIT Press, 2016.

    \item Dokumentacja TensorFlow/Keras:

          \url{https://www.tensorflow.org/api_docs/python/tf/keras}

    \item Dokumentacja scikit-learn (metryki klasyfikacji):

          \url{https://scikit-learn.org/stable/modules/model_evaluation.html}

\end{enumerate}



\end{document}

